% Nguồn phần đầu đã đối soát; phần chưa ký giữ nhãn rõ.
\chapter*{Acknowledgements}
\addcontentsline{toc}{chapter}{Acknowledgements}
\markboth{Acknowledgements}{Acknowledgements}
I thank my supervisor, Dr.~Tran Hong Viet.

\chapter*{Declaration}
\addcontentsline{toc}{chapter}{Declaration}
\markboth{Declaration}{Declaration}
\textbf{Unsigned declaration.} Author-confirmed wording; signature and date remain to be completed by the author.

I acknowledge the published methods, software and data sources on which
this thesis builds, and distinguish inherited components from the
adaptations and evaluation reported here. The historical findings are
presented with their adverse outcomes and unresolved limitations.
Unperformed experiments are identified as future work, and historical
simulations are distinguished from prospective trading evidence.

AI-assisted tools have been used to draft and edit this thesis manuscript.
Responsibility for checking the final text, references, attribution and
reported evidence rests with the author.

\textbf{Author:} Phan Hoang Dung\\
\textbf{Student ID:} 23020346\\
\textbf{Signature:} {[}To be completed by the author{]}\\
\textbf{Date:} {[}To be completed by the author{]}

\chapter*{Abstract}
\addcontentsline{toc}{chapter}{Abstract}
\markboth{Abstract}{Abstract}
Large language models are increasingly used in financial systems that
divide analysis across specialist agents. TradingAgents coordinates
market, news and fundamental analysis, QuantAgent assigns agents to
technical indicators, chart patterns and trends, and FinMem retrieves
information from layered financial memory. These approaches provide
different ways to describe a market or retrieve its history. For daily
equity decisions, an open question is whether completed trading episodes
from comparable market regimes can improve decisions, and when the
historical outcomes are too sparse to justify changing an action.

This thesis studies regime-aware Bayesian in-context learning (BICL)
on top of Our system, which adds quantitative alpha and Vietnamese
news sentiment to QuantAgent. BICL v1--v4 vary how eligible episodes
are matched, summarized and presented to the Decision Agent. Episodes
must have ended before the current decision. BICL v5 keeps past
trading tasks distinct from observations of the target task and uses a
numerical teacher--student model to estimate the probability of a
positive net outcome. It permits a revision to the fact-checked
reference action when a separate utility check supports the change.
Each variant retains its designated reference action when its
support rules require it.

The follow-up records 2,448 forecasts across eight methods, three
Decision repeats and 102 stock--dates. Each cutoff supplies one shared
upstream analysis. The execution rules yield 96 labels. The secondary
comparison of return and balanced accuracy uses 90 cutoffs in 45
accounts with both scheduled cycles per repeat. Full-cohort endpoints
therefore remain undefined. The secondary simultaneous bands are also
undefined because some contrast draws have zero variance.

BICL v1--v5 change final actions at two, one, one, one and one distinct
stock--dates, respectively. All three v5 revisions occur at the same
HVN cutoff, on 1 July 2024. Versions v1--v4 yield returns close to the
shared-prompt baseline, and v5 has lower mean return than both baselines
and Fact Anchor on the complete accounts. The follow-up measures the
effect of BICL under sparse historical support and frequent action reuse.

\textbf{Keywords:} Vietnamese stock market; multi-agent large language models; regime-aware
context; Bayesian in-context learning; LONG/CASH; point-in-time evaluation.
